\vfill\pagebreak
\section{Looking up by key: maps}
\label{sec:collections:maps}

An array finds a value by its position, a number from 0 to its length minus
one. That suits the months, which are numbered, but most tables are looked up by
something else. A dictionary gives the definition of a word, a phone book the
number of a person, a gradebook the grade of a student.

Stored in an array, a phone book would leave two ways to find the number of
Alice. The first is to know that she is the 312th entry, which nobody knows or
wants to know. The second is to go through the entries one by one until her
name comes up, which takes longer as the book grows.

A \textit{map} is a table looked up by something other than a position. It
associates \textit{keys} with values: the words of a dictionary with their
definitions, the names of a phone book with their numbers. Each key appears once
in a map, with one value, and the map finds the value of a key directly,
without going through the other entries.

\lstinputlisting[style=coloredverbatim, caption={Capitals by country, \textit{examples/collections/capitals.yr}}, label=lst:collections:capitals]{examples/collections/capitals.yr}
\vspace{-10pt}%
\begin{lstlisting}[style=bashVerb, escapechar=@+]
@+\textcolor{teal!80}{alice@dev:\mtilde}@+$ ./capitals
4 countries
Country: Peru
The capital of Peru is Lima.
@+\textcolor{teal!80}{alice@dev:\mtilde}@+$ ./capitals
4 countries
Country: Chile
I do not know Chile.
\end{lstlisting}

\subsection{Writing a map}

A map is written as its entries between square brackets, separated by commas.
An entry is a key and its value, separated by \token{=>}, as on lines~4 to~7.
By convention, the entries are written one per line, unless there are only two
of them, or a great many. The compiler does not enforce it.

The type of a map is written the same way, with the type of the keys and the
type of the values: \token{capitals} is a \token{[[c8] => [c8]]}, a map from
strings to strings. All the keys have the same type, and so do all the values,
but the two types may differ.

\begin{lstlisting}[style=coloredverbatim]
let ages = copy ["Alice" => 31, "Bruno" => 27];
let squares = copy [1 => 1,
                    2 => 4,
                    3 => 9];
println(ages, " ", squares);
\end{lstlisting}

\token{ages} is a \token{[[c8] => i32]}, from strings to integers, and
\token{squares} an \token{[i32 => i32]}, from integers to integers.

Like a slice, a map stores its entries away from the variable, and is made with
\token{copy} (Section~\ref{sec:collections:copy}). Its length is its number of
entries: \token{capitals.len}, on line~8, is 4.

\subsection{Looking up a key}

\token{capitals[country]}, on line~11, looks up the value of the key
\token{country}. The key may be missing, so the result is not a string, but an
option (Section~\ref{sec:collections:options}) that holds the string when the key
is there.

\begin{lstlisting}[style=coloredverbatim]
let capitals = copy ["France" => "Paris", "Japan" => "Tokyo"];
println(capitals["Japan"]); // Ok(Tokyo)
println(capitals["Chile"]); // Err()
\end{lstlisting}

\token{if let} takes the value out when there is one, as on line~11.

When the value itself is not needed, \token{key in m} tests whether the map
\token{m} has the key, and \token{key !in m} whether it has not. On a map,
\token{in} tests the keys, not the values.

\begin{lstlisting}[style=coloredverbatim]
let capitals = copy ["France" => "Paris", "Japan" => "Tokyo"];
println("Japan" in capitals);  // true
println("Chile" !in capitals); // true
println("Paris" in capitals);  // false, "Paris" is a value
\end{lstlisting}

\subsection{Adding and changing entries}

\token{m[key] = value} gives \token{key} the value \token{value}. If the map
already had the key, its old value is replaced; if not, the entry is added. The
entries are the elements of the map, so the rule of
Section~\ref{sec:collections:dmut} applies, and the map must be declared
\token{dmut}.

\begin{lstlisting}[style=coloredverbatimError, escapechar=@, caption=Adding an entry to a map declared with \token{let}]
use std::io;

fn main() {
  let ages = copy ["Alice" => 31, "Bruno" => 27];
  @\hb{ages["Chloe"]}@ = 22; // error
  println(ages);
}
\end{lstlisting}

The compiler says \errmsg{E4151}{not a lvalue}: \token{ages["Chloe"]} cannot
stand on the left of an assignment. Declared with \token{let dmut ages}, the
program is correct.

A map is often built from nothing, one entry at a time. The next listing counts
how many times each letter appears in a word. It starts from an empty map,
\token{copy []}, whose type must then be written, since there is no entry to
infer it from: \token{[c8 => i32]}, from characters to counts.

\lstinputlisting[style=coloredverbatim, caption={Counting the letters of a word, \textit{examples/collections/letters.yr}}, label=lst:collections:letters]{examples/collections/letters.yr}
\vspace{-10pt}%
\begin{lstlisting}[style=bashVerb, escapechar=@+]
@+\textcolor{teal!80}{alice@dev:\mtilde}@+$ ./letters
Word: mississippi
p: 2
i: 4
m: 1
s: 4
\end{lstlisting}

For each letter, the \token{match} of lines~7 to~10 looks the letter up. A letter
already seen has a count, \token{n}, and the arm on line~8 replaces it with
\token{n + 1}. A letter seen for the first time has none, and the arm on line~9
adds it, with a count of 1. An arm that is a single statement, such as these
assignments, can be written without braces, ended by its semicolon.

\notebox{An empty slice is written \token{[]}, but an empty map needs
  \token{copy []}. Without the \token{copy}, the declaration of line~5 is
  refused with \errmsg{E4095}{incompatible types mut [c8 => mut i32] and mut
    [void ; 0us]}, and the note under it says that creating a map on the stack
  is prohibited. Why a slice can do without a copy when it is empty, and a map
  cannot, depends on where each keeps its elements. Chapter~\ref{chap:layout}
  explains it.}

\subsection{Going through a map}

\token{for} goes through the entries of a map, as it goes through the elements
of a slice. The loop of line~12 of Listing~\ref{lst:collections:letters} gives
it two names, \token{letter} and \token{count}: the first holds the key of the
current entry, and the second its value. The next listing shows the other forms.

\lstinputlisting[style=coloredverbatim, caption={Going through a map, \textit{examples/collections/prices.yr}}, label=lst:collections:prices]{examples/collections/prices.yr}
\vspace{-10pt}%
\begin{lstlisting}[style=bashVerb]
cheese costs 7
bread costs 2
milk costs 1
cheese bread milk
Total: 10
[cheese => 14, bread => 4, milk => 2]
\end{lstlisting}

\begin{itemize}
\item With two names, on line~8, the loop goes through the keys and their
  values.
\item With a single name, on line~12, it goes through the keys only.
\item With the wildcard in place of the key, on line~18, it goes through the
  values only.
\end{itemize}

As for arrays, the names are constants. \token{price = price * 2;} in the loop
of line~23 would be refused with \errmsg{E4153}{price is a value iterator, and
  cannot be used as a lvalue}: \token{price} is a copy of the value, and changing
it would not change the map. A value is changed through the map itself, with
\token{prices[item] = price * 2;}, as on line~24. That requires the map to be
declared \token{dmut}, as it is on line~4.

\subsection{The order of the entries}

The entries do not come out in the order they were written in: \token{bread}
was the first, and the loops of Listing~\ref{lst:collections:prices} list it
second. For the letters of \textit{mississippi}, the order was neither that of
the word, nor the alphabetical order: the \textit{m} came third.

A map arranges its entries in the order that makes finding a key fast. That
order depends on the keys, and on how many entries the map holds, and a program
cannot rely on it. When the order matters, the program goes through the keys it
wants, in the order it wants, and looks each one up.

\begin{lstlisting}[style=coloredverbatim]
let ages = copy ["Alice" => 31,
                 "Bruno" => 27,
                 "Chloe" => 22];
for name, age in ages {
  println(name, " ", age);
}
for name in ["Alice", "Bruno", "Chloe"] {
  if let Ok(age) = ages[name] {
    println(name, " ", age);
  }
}
\end{lstlisting}
\vspace{-10pt}%
\begin{lstlisting}[style=bashVerb]
Chloe 22
Alice 31
Bruno 27
Alice 31
Bruno 27
Chloe 22
\end{lstlisting}

The first loop follows the order of the map, the second that of the array of
names. The program of Section~\ref{sec:collections:gradebook} uses the same
method.

\subsection{Removing an entry}

An entry is removed with \token{m:.remove(key)}. The colon in front of the dot
marks a call that changes \token{m} itself, so the map must be declared
\token{dmut}. A later chapter explains that colon; until then, it is a formula.

\begin{lstlisting}[style=coloredverbatim]
let dmut ages = copy ["Alice" => 31,
                      "Bruno" => 27,
                      "Chloe" => 22];
ages:.remove("Bruno");
println(ages.len, " ", "Bruno" in ages); // 2 false
ages:.remove("Zoe");
println(ages.len);                       // 2
\end{lstlisting}

Removing a key that the map does not have, such as \token{"Zoe"}, does nothing.

\subsection{Sets}

Some tables have keys and nothing to associate with them: the words already seen
in a text, the numbers already drawn in a lottery. Such a collection, where each
value appears once and the only question is whether a value is in it, is a
\textit{set}. Ymir has no set type of its own: a set is a map whose values are
the empty tuple, \token{()}, a tuple of no field, which holds no information. A
set of strings is a \token{[[c8] => ()]}.

\begin{lstlisting}[style=coloredverbatim]
let dmut seen: [[c8] => ()] = copy [];
for word in ["the", "cat", "saw", "the", "dog"] {
  seen[word] = ();
}
println(seen.len, " ", "cat" in seen, " ", "cow" in seen);
seen:.remove("cat");
for word in seen {
  print(word, " ");
}
println();
\end{lstlisting}
\vspace{-10pt}%
\begin{lstlisting}[style=bashVerb]
4 true false
the dog saw
\end{lstlisting}

\token{seen[word] = ();} adds a word to the set. A word added twice, such as
\token{"the"}, is still there once, so the set holds 4 words, not 5.
\token{in} tests whether a word is in the set, \token{:.remove} takes one out,
and a \token{for} with a single name goes through the words, in the order of
the map.
